%%%%%%%%%%%%%%%%%%%%%%%%%%%%%%%%%%%%%%%%%%%%%%%%%%%%%%%%%%%%%%%%%%%%%%
% Lớp 10 | Bất phương trình bậc 2 một ẩn | Phương trình quy về phương trình bậc hai
% Dạng 0D7?3-4: Phương trình căn thức (dạng khác)
% Số câu: 5  (sinh tự động bởi tach_cau.py; từ nay thầy sửa trực tiếp file này)
%%%%%%%%%%%%%%%%%%%%%%%%%%%%%%%%%%%%%%%%%%%%%%%%%%%%%%%%%%%%%%%%%%%%%%

%%%%%%%%%%%%%%%% TRẢ LỜI NGẮN %%%%%%%%%%%%%%%%

% nguồn: MaHoaNH_TraLoiNgan_Lop10
\begin{bt}%[0D7H3-4]%[ma:df484765]
	Phương trình $\sqrt{x-4}\cdot\left(x^2-3x+2\right)=0$  có bao nhiêu nghiệm?
	 \shortans{1} \loigiai{
	Điều  kiện $x-4\geqslant 0 \Leftrightarrow x\geqslant4$. \\
	Phương trình đã cho tương  đương với
	$$\hoac{ & \sqrt{x-4}=0 \\ & x^2-3x+2=0} \Leftrightarrow \hoac{ & x=4 & & \text{(thỏa mãn)} \\ & x=2 & & \text{(loại)}  \\&x=1 & & \text{(loại)}.}$$
	Vậy phương trình  $\sqrt{x-4}\cdot\left(x^2-3x+2\right)=0$ có $1$ nghiệm là $x=4$.
	} 
\end{bt}

% nguồn: MaHoaNH_TraLoiNgan_Lop10
\begin{bt}%[0D7C3-4]%[ma:3c92b7e4]
	Tìm số nghiệm  của phương trình $\sqrt{x + 2\sqrt{x-1}} = x+\dfrac{1}{4}$.
	\shortans{1} 
	 \loigiai{
		Ta có $x+2\sqrt{x-1} = x-1 + 2\sqrt{x-1} + 1 = \left( \sqrt{x-1} + 1 \right)^2$ và $\sqrt{x-1} + 1 >  0$ với mọi $x\ge 1$ nên phương trình đã cho tương đương với
		\begin{eqnarray*}
			& &\sqrt{x-1} + 1 = x +  \dfrac{1}{4}
			\Leftrightarrow  4\sqrt{x-1} = 4x-3 \\
			&\Leftrightarrow  &\heva{&4x-3\ge 0\\ &16(x-1) = (4x-3)^2}
			\Leftrightarrow \heva{&x\ge \dfrac{3}{4}\\ &16x^2 - 40x + 25 =  0}\\
			&\Leftrightarrow & \heva{&x\ge \dfrac{3}{4}\\  & x = \dfrac{5}{4}} \Leftrightarrow x = \dfrac{5}{4}.
		\end{eqnarray*}
		 Vậy phương trình có nghiệm duy nhất.
	 }
\end{bt}

% nguồn: MaHoaNH_TraLoiNgan_Lop10
\begin{bt}%[0D7C3-4]%[ma:c9f9aaf9]
	Tính tổng bình phương các nghiệm của  phương trình
	 \[
	x^2 - 3x + 10 +  2(x-3)\sqrt{3x+1} = 0.
	 \]
	 \shortans{1}
	 \loigiai{
		Ta có phương  trình tương đương với
		\begin{eqnarray*}
			&  &\Leftrightarrow (x^2 - 6x + 9) + (3x+1) + 2(x-3)\sqrt{3x+1} = 0 \\
			& &\Leftrightarrow (x-3)^2 + \left( \sqrt{3x+1} \right)^2 + 2(x-3)\sqrt{3x+1} = 0\\ 
			& &\Leftrightarrow \left( x - 3 + \sqrt{3x+1} \right) ^ 2 = 0 \Leftrightarrow x - 3 + \sqrt{3x+1} = 0 \Leftrightarrow  \sqrt{3x+1} = 3-x\\
			&  & \Leftrightarrow \heva{&3-x\ge 0\\ &3x+1 = (3-x)^2}
			\Leftrightarrow \heva{&x\le 3\\ &x^2 - 9x + 8 =  0}
			\Leftrightarrow  \heva{&x\le 3\\ &\hoac{&x=1\\ &x=8}}
			\Leftrightarrow  x = 1.
		\end{eqnarray*}
		Vậy tổng bình phương tất cả các nghiệm của phương trình là $1$. 
	 }
\end{bt}

% nguồn: MaHoaNH_TraLoiNgan_Lop10
\begin{bt}%[0D7C3-4]%[ma:c9025f52]
	Biết rằng phương trình $(x+1)(x+4) - 5\sqrt{x^2+5x+2} = 6$ có hai nghiệm phân biệt  $a$ và $b$ thoả mãn $a < b$. Tính giá trị biểu thức $T = 2a + 12b$.
	\shortans{10} 
	 \loigiai{
		Điều kiện xác  định $x^2+5x+2\ge 0$.\\
		Ta có 
		 \[
		 (x+1)(x+4) - 5\sqrt{x^2+5x+2} = 6 \Leftrightarrow x^2 + 5x - 2 - 3\sqrt{x^2 + 5x+2} = 0.
		 \]
		Đặt $t = \sqrt{x^2 + 5x+2}$ ($t\ge 0$) thì $x^2+5x -2 = t^2 - 4$. Khi  đó phương trình trở thành
		\[ 
		t^2 - 3t - 4 = 0 \Leftrightarrow \hoac{&t = 4 &  & \text{ (thoả mãn)}\\ &t = -1 & & \text{ (loại)}.}
		\] 
		Với  $t=4$ thì
		 \[
		\sqrt{x^2+5x+2} = 4\Leftrightarrow x^2 + 5x +2 = 16 \Leftrightarrow x^2 + 5x - 14 = 0 \Leftrightarrow 
		\hoac{&x  = 2\\ &x=-7.}
		\] 
		Vậy $a=-7$ và $b=2$ nên $T  = 2a+14b = 10$.
	 }
\end{bt}

% nguồn: MaHoaNH_TraLoiNgan_Lop10
\begin{bt}%[0D7C3-4]%[ma:032a649d]
	Tính tổng các nghiệm của phương trình $x^2 -3x + 1 +\sqrt{2x-1}  = 0$.
	 \shortans{1,59}
	 \loigiai{
		Điều kiện xác định $x\ge \dfrac{1}{2}$. Khi đó phương trình tương  đương
		\begin{eqnarray*}
			&  &\sqrt{2x-1} - x + x^2 - (2x-1) = 0 \Leftrightarrow
			\left( \sqrt{2x-1} - x \right) + \left( x - \sqrt{2x-1}\right)\left( x +  \sqrt{2x-1} \right) = 0 \\
			&\Leftrightarrow & \left( x - \sqrt{2x-1} \right)\left( -1 + x +  \sqrt{2x-1} \right)
			= 0  \Leftrightarrow \hoac{&\sqrt{2x-1} = x\\ &\sqrt{2x-1} = 1-x}\\
			 &\Leftrightarrow
			&\hoac{ 
			&\heva{&x\ge 0\\  & 2x-1 = x^2} &\Leftrightarrow \heva{&x\ge 0\\ &x^2 - 2x + 1 =0}
			&\Leftrightarrow x = 1\\ 
			 &\heva{&1-x\ge 0\\ & 2x-1 = (1-x)^2} &\Leftrightarrow \heva{&x\le 1\\ &x^2 - 4x + 2 =0}
			&\Leftrightarrow x = 2 -  \sqrt{2}.\\
			 }
		\end{eqnarray*}
		Vậy phương  trình có tập nghiệm là $S = \left\{1;2-\sqrt{2}\right\}$ nên tổng các nghiệm của phương trình là
		\[ 
		1 + 2-\sqrt{2  } = 3-\sqrt{2}.
		\] 
	 }
\end{bt}
